\chapter{1977年江苏无锡市初试卷}

\section{填空题}

\begin{problem}
\begin{enumerate}
\item 求代数式 \(f(x)=3x^2-2x-5\) 的值，当 \(x=0\) 时，\(f(0)=\fillinblank{}\)；当 \(x=-1\) 时，\(f(-1)=\fillinblank{}\)；当 \(x=\frac{1}{3}\) 时，\(f\left(\frac{1}{3}\right)=\fillinblank{}\).
\item 在直角三角形 \(ABC\) 中，已知 \(\angle C=90^\circ\)，\(AB=13\)，\(BC=5\)，则 \(\sin A=\fillinblank{}\)，\(\cos B=\fillinblank{}\)，\(\tan A=\fillinblank{}\)，\(\cot A=\fillinblank{}\).
\item 化简：\(\sqrt{(3-x)^2}\).
\item 已知圆方程为 \(x^2+y^2=25\)，那么这个圆以原点为极点，\(x\) 轴的正方向为极轴的极坐标方程是\fillinblank{}，圆心的极坐标是\fillinblank{}.
\item 已知长方体的长，宽，高分别是 \(3\mathrm{~cm}\)，\(2\mathrm{~cm}\)，\(5\mathrm{~cm}\)，
则它的体积为\fillinblank{}，它的表面积为\fillinblank{}.
\end{enumerate}
\end{problem}

\begin{answer}
\begin{enumerate}
\item \(-5\)，\(0\)，\(-\frac{16}{3}\).
\item \(\frac{5}{13}\)，\(\frac{5}{13}\)，\(\frac{5}{12}\)，\(\frac{12}{5}\).
\item \(\left|3-x\right|\).
\item \(r=5\)，\((0,0)\).
\item \(30\mathrm{~cm}^3\)，\(62\mathrm{~cm}^2\).
\end{enumerate}
\end{answer}

\begin{solution}
\begin{enumerate}
\item 当 \(x=0\) 时，\(f(0)=3\times0^2-2\times0-5=-5\)；当 \(x=-1\) 时，\(f(-1)=3\times(-1)^2-2\times(-1)-5=3+2-5=0\)；当 \(x=\frac{1}{3}\) 时，\(f\left(\frac{1}{3}\right)=3\times\left(\frac{1}{3}\right)^2-2\times\frac{1}{3}-5=\frac{1}{3}-\frac{2}{3}-5=-\frac{16}{3}\).
\item 因为 \(\triangle ABC\) 在 \(C\) 处直角，所以 \(AB\) 是斜边，\(AC=\sqrt{AB^2-BC^2}=\sqrt{13^2-5^2}=\sqrt{144}=12\). 因而 \(\sin A=\frac{BC}{AB}=\frac{5}{13}\)，\(\cos B=\frac{BC}{AB}=\frac{5}{13}\)，\(\tan A=\frac{BC}{AC}=\frac{5}{12}\)，\(\cot A=\frac{AC}{BC}=\frac{12}{5}\).
\item 对任意实数 \(x\)，有
\[
\sqrt{(3-x)^2}=\left|3-x\right|=\begin{cases}
3-x,&x\leq3,\\
x-3,&x>3
\end{cases}
\]
所以化简结果为 \(\left|3-x\right|\).
\item 令极坐标为 \((r,\varphi)\)，则 \(x^2+y^2=r^2\)，原方程化为 \(r^2=25\). 由于极径 \(r\geq0\)，故极坐标方程为 \(r=5\). 该圆的圆心就是原点，也就是极点，因此圆心的极坐标可记为 \((0,0)\)（极角任意）.
\item 长方体的体积为 \(3\times2\times5=30\mathrm{~cm}^3\)，表面积为 \(2\left(3\times2+2\times5+3\times5\right)=2\times31=62\mathrm{~cm}^2\).
\end{enumerate}
\end{solution}

\section{解答题}

\begin{problem}
计算：

\begin{enumerate}
\item \(\left[\left|8-6\right|-\left|4-7\right|\right]\times\frac{1}{\left|1-2\right|}\).
\item \(a^{\frac{1}{2}}b^0\left(ab^2\right)^{-\frac{1}{2}}\div\left[b^{\frac{1}{3}}\left(a^3b^{-2}\right)^{-\frac{1}{3}}\right]\).
\item \(\left(\lg50\right)^2+\lg2\cdot\lg\left(50^2\right)+\left(\lg2\right)^2\).
\item \(\sin105^\circ+\sin15^\circ\).
\item \(\displaystyle\lim_{n\to\infty}\frac{5-2n+3n^2}{4n^2-n-7}\).
\item \(\sqrt{\frac{1-\cos\theta}{1+\cos\theta}}+\sqrt{\frac{1+\cos\theta}{1-\cos\theta}}\).
\end{enumerate}
\end{problem}

\begin{solution}
\begin{enumerate}
\item \(\left|8-6\right|=2\)，\(\left|4-7\right|=3\)，\(\left|1-2\right|=1\)，所以原式 \(=(2-3)\times\frac{1}{1}=-1\).
\item 按有理指数幂的通常约定取 \(a>0\)，\(b>0\). 因为 \(b^0=1\)，且
\[
a^{\frac12}b^0\left(ab^2\right)^{-\frac12}=a^{\frac12}\cdot1\cdot a^{-\frac12}b^{-1}=b^{-1}
\]
又
\[
b^{\frac13}\left(a^3b^{-2}\right)^{-\frac13}=b^{\frac13}a^{-1}b^{\frac23}=a^{-1}b
\]
所以原式 \(=b^{-1}\div\left(a^{-1}b\right)=\frac{1}{b}\times\frac{a}{b}=\frac{a}{b^2}\).
\item 设 \(x=\lg2\)，\(y=\lg50\). 由对数的乘方和积的性质，\(\lg\left(50^2\right)=2\lg50=2y\)，且 \(x+y=\lg2+\lg50=\lg100=2\). 因此原式
\[
y^2+x\cdot2y+x^2=(x+y)^2=2^2=4
\]
\item 由和差化积公式，\(\sin105^\circ+\sin15^\circ=2\sin60^\circ\cos45^\circ=2\times\frac{\sqrt3}{2}\times\frac{\sqrt2}{2}=\frac{\sqrt6}{2}\).
\item 分子、分母同除以 \(n^2\)，得
\[
\displaystyle\lim_{n\to\infty}\frac{5-2n+3n^2}{4n^2-n-7}=\lim_{n\to\infty}\frac{\frac{5}{n^2}-\frac{2}{n}+3}{4-\frac{1}{n}-\frac{7}{n^2}}=\frac{3}{4}
\]
\item 题式有意义时，需有 \(1-\cos\theta\neq0\) 且 \(1+\cos\theta\neq0\)，即 \(\sin\theta\neq0\). 利用 \(\sin^2\theta=(1-\cos\theta)(1+\cos\theta)\)，有
\[
\sqrt{\frac{1-\cos\theta}{1+\cos\theta}}=\sqrt{\frac{(1-\cos\theta)^2}{\sin^2\theta}}=\frac{1-\cos\theta}{\left|\sin\theta\right|}
\]
\[
\sqrt{\frac{1+\cos\theta}{1-\cos\theta}}=\sqrt{\frac{(1+\cos\theta)^2}{\sin^2\theta}}=\frac{1+\cos\theta}{\left|\sin\theta\right|}
\]
两式相加，原式 \(=\frac{2}{\left|\sin\theta\right|}\).
\end{enumerate}
\end{solution}

\begin{problem}
解方程、方程组：

\begin{enumerate}
\item \(\frac{5}{x-2}+\frac{2}{x^2+2x}=\frac{20}{x^2-4}\).
\item \(\begin{cases}
x-\frac{y-1}{2}=3,\\
\frac{x}{4}+y=4
\end{cases}
\)
\item 当 \(m\) 为何值时，二次方程 \(mx^2+3mx+3=0\)

\begin{enumerate}
\item 有两个相等的实数根；
\item 无实数根.
\end{enumerate}
\end{enumerate}
\end{problem}

\begin{solution}
\begin{enumerate}
\item 原方程有意义的条件是 \(x\neq2\)，\(x\neq0\)，\(x\neq-2\). 在此条件下，两边同乘 \(x(x-2)(x+2)\)，得
\[
5x(x+2)+2(x-2)=20x
\]
即 \(5x^2-8x-4=0\). 因为
\[
5x^2-8x-4=(5x+2)(x-2)
\]
所以 \(x=-\frac25\) 或 \(x=2\). 其中 \(x=2\) 使原方程分母为零，必须舍去，故原方程的解为 \(x=-\frac25\).
\item 第一个方程两边乘以 \(2\)，得 \(2x-y+1=6\)，所以 \(y=2x-5\). 代入第二个方程，得 \(\frac{x}{4}+2x-5=4\)，即 \(\frac{9x}{4}=9\)，从而 \(x=4\)，再代回得 \(y=3\). 所以方程组的解为 \(\begin{cases}
x=4,\\
y=3
\end{cases}\).
\item 因为题目中的方程是二次方程，所以先要求 \(m\neq0\). 它的判别式为 \(\Delta=(3m)^2-4m\cdot3=3m(3m-4)\).
\begin{enumerate}
\item 有两个相等的实数根的条件是 \(\Delta=0\). 由 \(3m(3m-4)=0\) 得 \(m=0\) 或 \(m=\frac43\)，但 \(m=0\) 时原式变为 \(3=0\)，不是二次方程，故 \(m=\frac43\).
\item 无实数根的条件是 \(\Delta<0\)，即 \(3m(3m-4)<0\). 在 \(m\neq0\) 的条件下，解得 \(0<m<\frac43\).
\end{enumerate}
\end{enumerate}
\end{solution}

\begin{problem}
分别写出数列 \(1\)，\(3\)，\(5\)，\(7\)，\(\cdots\) 和 \(1\)，\(\frac{1}{2}\)，\(\frac{1}{4}\)，\(\frac{1}{8}\)，\(\cdots\) 的通项公式，并用求和公式计算：\(1+1\)，\(3+\frac{1}{2}\)，\(5+\frac{1}{4}\)，\(7+\frac{1}{8}\)，\(\cdots\)，前10项的和.
\end{problem}

\begin{solution}
设第一个数列为 \(\{a_n\}\)，第二个数列为 \(\{b_n\}\). 第一个数列是首项为 \(1\)、公差为 \(2\) 的等差数列，所以
\[
a_n=1+(n-1)\times2=2n-1
\]
第二个数列是首项为 \(1\)、公比为 \(\frac12\) 的等比数列，所以
\[
b_n=1\times\left(\frac12\right)^{n-1}=\left(\frac12\right)^{n-1}
\]
题中各项分别为 \(a_n+b_n\)，即通项为 \(2n-1+\left(\frac12\right)^{n-1}\). 因而前 \(10\) 项的和为
\[
S_{10}=\sum_{n=1}^{10}\left(2n-1\right)+\sum_{n=1}^{10}\left(\frac12\right)^{n-1}
=\frac{10(1+19)}{2}+\frac{1-\left(\frac12\right)^{10}}{1-\frac12}
\]
\[
S_{10}=100+\frac{1023}{512}=\frac{52223}{512}
\]
所以前 \(10\) 项的和为 \(\frac{52223}{512}\).
\end{solution}

\begin{problem}
在 \(\triangle ABC\) 中，\(\angle A=60^\circ\)，\(AB=2\)，\(AC=1+\sqrt{3}\)，解这个三角形.
\end{problem}

\begin{solution}
在 \(\triangle ABC\) 中，已知 \(\angle A=60^\circ\)，\(AB=2\)，\(AC=1+\sqrt3\). 由余弦定理，
\[
BC^2=AB^2+AC^2-2\cdot AB\cdot AC\cos A=4+(1+\sqrt3)^2-2\cdot2\cdot(1+\sqrt3)\cdot\frac12=6
\]
故 \(BC=\sqrt6\).

再由余弦定理求 \(\angle B\)，
\[
\cos B=\frac{AB^2+BC^2-AC^2}{2\cdot AB\cdot BC}=\frac{4+6-(1+\sqrt3)^2}{4\sqrt6}=\frac{\sqrt6-\sqrt2}{4}=\cos75^\circ
\]
因为 \(0^\circ<B<180^\circ\)，所以 \(B=75^\circ\)，从而 \(C=180^\circ-60^\circ-75^\circ=45^\circ\).

因此，三角形的三边为 \(AB=2\)，\(AC=1+\sqrt3\)，\(BC=\sqrt6\)，三个内角为 \(A=60^\circ\)，\(B=75^\circ\)，\(C=45^\circ\).
\end{solution}

\begin{problem}
把底面半径和高均为 \(a\) 的一个圆柱体和圆锥体的金属物体，熔铸成一个金属球，求这个球的半径（损耗不计）.
\end{problem}

\begin{solution}
设熔铸成的金属球半径为 \(r\). 原圆柱体的体积为 \(\pi a^2\cdot a=\pi a^3\)，原圆锥体的体积为 \(\frac13\pi a^2\cdot a=\frac13\pi a^3\). 因为熔铸过程中没有损耗，所以金属球的体积满足
\[
\frac43\pi r^3=\pi a^3+\frac13\pi a^3=\frac43\pi a^3
\]
约去 \(\frac43\pi\)，得 \(r^3=a^3\). 半径为正数，故 \(r=a\).
\end{solution}

\begin{problem}
如果直线的倾角 \(\alpha\) 满足 \(4\sin\alpha=3\cos\alpha\)，并过椭圆 \(\frac{x^2}{16}+\frac{y^2}{25}=1\) 的焦点，求直线方程以及直线与坐标轴围成的三角形的面积.
\end{problem}

\begin{solution}
椭圆 \(\frac{x^2}{16}+\frac{y^2}{25}=1\) 的长半轴为 \(5\)，短半轴为 \(4\)，所以焦半距
\[
c=\sqrt{5^2-4^2}=3
\]
两个焦点为 \(F_1(0,3)\)，\(F_2(0,-3)\).

由直线倾角的条件 \(4\sin\alpha=3\cos\alpha\)，且直线倾角满足 \(0\leq\alpha<180^\circ\)，可知斜率
\[
k=\tan\alpha=\frac34
\]
分别过两个焦点作斜率为 \(\frac34\) 的直线，得到
\[
y=\frac34x+3\quad\text{或}\quad y=\frac34x-3
\]

对直线 \(y=\frac34x+3\)，令 \(y=0\)，得横截距为 \(-4\)；令 \(x=0\)，得纵截距为 \(3\)，所以它与坐标轴围成的三角形面积为 \(\frac12\times4\times3=6\). 对直线 \(y=\frac34x-3\)，横、纵截距分别为 \(4\)，\(-3\)，所围三角形面积同样为 \(\frac12\times4\times3=6\).
\end{solution}

\begin{problem}
在 \(\odot O\) 中，直径 \(AB\) 和弦 \(AC\) 的夹角为 \(30^\circ\)，切线 \(CD\) 和 \(AB\) 的延长线交于 \(D\)，求证：\(AC^2=AO\cdot AD\).

\FigureLayoutDeclare{img/1977/jiangsu_wuxi_city/q8_fig1.png}{6.0cm}{24bp 0bp 0bp 0bp}
\bitmapfigure[max width=5.2cm]{img/1977/jiangsu_wuxi_city/q8_fig1.png}
\end{problem}

\begin{solution}
因为 \(AB\) 是直径，所以 \(\angle ACB=90^\circ\). 又已知直径 \(AB\) 与弦 \(AC\) 的夹角为 \(\angle BAC=30^\circ\)，故在直角三角形 \(ABC\) 中，\(\angle ABC=60^\circ\). 于是
\[
AC=AB\cos30^\circ=2AO\cdot\frac{\sqrt3}{2}=\sqrt3\,AO
\]

由切线与弦所夹角的定理，弦 \(AC\) 与切线在 \(C\) 处所成的锐角等于 \(\angle ABC=60^\circ\). 因为 \(D\) 在切线的另一侧射线上，三角形 \(ACD\) 的内角 \(\angle ACD=180^\circ-60^\circ=120^\circ\). 又 \(\angle CAD=30^\circ\)，所以 \(\angle ADC=30^\circ\). 在 \(\triangle ACD\) 中应用正弦定理，
\[
\frac{AD}{\sin120^\circ}=\frac{AC}{\sin30^\circ}
\]
从而 \(AD=\sqrt3\,AC\). 因此
\[
AO\cdot AD=AO\cdot\sqrt3\,AC=(\sqrt3\,AO)\,AC=AC^2
\]
所以得证.
\end{solution}

\begin{problem}
证明：如果两数的和为一定值，那么当两数相等时，它们的乘积最大.
\end{problem}

\begin{solution}
设两个数分别为 \(x\)，\(y\)，且它们的和为定值 \(s\)，即 \(x+y=s\). 则
\[
xy=\frac{(x+y)^2-(x-y)^2}{4}=\frac{s^2}{4}-\frac{(x-y)^2}{4}\leq\frac{s^2}{4}
\]
因为 \((x-y)^2\geq0\)，等号成立当且仅当 \(x-y=0\)，即 \(x=y=\frac{s}{2}\). 所以当两数相等时，它们的乘积最大.
\end{solution}
